\scanpage{33}

On introduit aussi une application (\uline{pas} un homomorphisme !)

\[
\left\{\begin{aligned}
 \deleted{\unsure{v}}\sigma_0 &:L_0^*\longrightarrow H\\
 \sigma_0(\beta)&\deleted{\longmapsto}=\bigl(\begin{smallmatrix}0&\deleted{\unsure{-}}\beta^{\deleted{\unsure{-1}}}\\-\beta^{-1}&0\end{smallmatrix}\bigr)
\end{aligned}\right.\tag*{(19)}
\]

\margin{\uline{NB} Si $L_0$ n'est pas libre, alors $L_0^*$ est vide, et ce qui suit sans intérêt. On suppose donc par la suite $L_0$ \deleted{\unsure{dual|trivial}} \inserted{libre}, donc $L_1$ libre}

on a alors

\[
\left\{\begin{aligned}
 \sigma_0(\beta)^2&=\bigl(\begin{smallmatrix}-1&0\\0&-1\end{smallmatrix}\bigr)=u(-1)\overset{\mathrm{def}}{=}\unsure{\omega|w}
 \quad\text{(élément \uline{central} de $H$)}\\
 \sigma_0(\beta)(u(\lambda))&=u(\lambda^{-1})\quad\text{donc }v(\beta)\text{ normalise }T\\
 \sigma_0(\lambda\beta\deleted{\unsure{\lambda|*}})&=u(\lambda)\sigma_0(\beta)
 =\underset{\deleted{\text{\unsure{donc pour}}}}{\sigma_0(\beta)u(\lambda^{-1})},
\end{aligned}\right.\tag*{(20)}
\]

la dernière formule montre que, si on pose

\[
\deleted{\unsure{L_0}}N^-=v_0(L_0^*)=\deleted{H}\left\{\bigl(\begin{smallmatrix}a&b\\c&d\end{smallmatrix}\bigr)\in H\mid a=0,d=0\right\},\tag*{(21)}
\]

alors $N^-$ est une classe à gauche ou à droite mod $T$, et que

\[
L_0^*\xrightarrow{\sim}N^-\tag*{(22)}
\]

est bijection, et \={m} un isomorphisme de $\Lambda^*$-torseurs, quand on fait opérer $\Lambda^*$ sur $L_0^*$ par homothéties, et sur $N^-$ par multiplication \uline{à gauche} par $u(\lambda)$. On pourrait bien sûr introduire \deleted{aussi le composé} \inserted{symétriquement}

\[
\sigma_1:L_1\xrightarrow{\sim}N^-\qquad \sigma_1(\gamma)=\bigl(\begin{smallmatrix}0&-\gamma^{-1}\\\gamma&0\end{smallmatrix}\bigr)\tag*{\deleted{(23)} (23)}
\]

de sorte que l'on a

\[
\left\{\begin{aligned}
 \sigma_1(\gamma)&=\sigma_0(-\gamma^{-1})\\
 \sigma_1(\lambda\gamma)&=u(\lambda^{-1})\sigma_1(\gamma)=\sigma_1(\gamma)u(\lambda)
\end{aligned}\right.\tag*{(24)}
\]

donc $\sigma_1$ est un isomorphisme de \inserted{$\Lambda^*$-}torseurs, quand cette fois on fait opérer $\Lambda^*$ sur $N^-$ par multiplication à \uline{droite} par $u(\lambda)$.

% archive sheet number: 32
